\section{Original Scheduling, Cadence, and Latency}
\label{sec:scheduling}
Let $H\ge1$ be the requested computation horizon in sample calls and let $B>0$
be the number of inner butterflies. The implementation chooses
\begin{equation}
 q=\ceil{B/H}
 \label{eq:quota}
\end{equation}
and invokes the single-step method $q$ times per call. Once complete, further
step invocations are no-ops. The following statements concern this integer
model. These reusable forward classes compute the ceiling through single-precision
floating-point division; the reported experiments verify the model over the
tested range. Extremely large sizes or horizons require a separate analysis of
rounding and integer overflow. Zero horizons and unsupported transform lengths
are outside the forward-path contract examined here. We use powers of two
$N\ge4$ for the real transform.

\begin{proposition}[Completion bound]
The transform completes after exactly $L=\ceil{B/q}$ calls, with $L\le H$.
\end{proposition}
\begin{proof}
After $\ell$ calls the number of completed butterflies is
$\min(B,\ell q)$. The least $\ell$ reaching $B$ is $\ceil{B/q}$.
Since $qH\ge B$, this value cannot exceed $H$.
\end{proof}

This is a bound on operation counts, not execution time. In particular,
completion \emph{within} $H$ calls does not imply completion \emph{at} call $H$.
In the original snapshot, both visualization modules test completion at the start of their processing
routine, publish the previous result, buffer a new frame, and then advance the
new transform. In steady state, with constant parameters and active processing,
frame starts are therefore separated by $L$ samples rather than necessarily by
$H$ samples. Initialization, freeze, reset, and reconfiguration are separate
state transitions and are not covered by this steady-state statement.

\begin{table}[tb]
\centering
\begin{tabular}{@{}rrrrrr@{}}
\toprule
$N$ & $H=N/2$ & $B_r$ & $q$ & $L$ & $H/L$ \\
\midrule
1,024 & 512 & 2,304 & 5 & 461 & 1.111 \\
2,048 & 1,024 & 5,120 & 5 & 1,024 & 1.000 \\
4,096 & 2,048 & 11,264 & 6 & 1,878 & 1.091 \\
16,384 & 8,192 & 53,248 & 7 & 7,607 & 1.077 \\
\bottomrule
\end{tabular}
\caption{Derived real-FFT schedules. $H/L$ is the actual frame rate divided by the nominal rate when a new frame starts immediately after completion. These are operation counts, not timing measurements.}
\label{tab:cadence}
\end{table}


For $N=4096$ and $H=2048$, $q=6$ and $L=1878$. At 48 kHz the nominal
interval is 42.67 ms, whereas the original steady-state interval is
39.125 ms. The rate difference is about 9.1\%. This is not a numerical FFT
error; it is a difference between the parameter's nominal interpretation and
the scheduler's completion-driven frame clock.

\subsection{Timestamp semantics}
Suppose a complete input window is captured on call $t_r$ and its first batch
runs in that same call. Its final batch runs at $t_r+L-1$, and the module
publishes the result at the start of call $t_r+L$. Measured from the newest
sample, publication delay is consequently $L/f_s$. Relative to the midpoint
of the captured sample support, its age is
\begin{equation}
 A_{\mathrm{mid}}=\frac{(N-1)/2+L}{f_s}.
 \label{eq:age}
\end{equation}
The midpoint is a timestamp convention, not a universal group-delay claim for
arbitrary windows. The display thread may add further delay. Figure~\ref{fig:timeline}
separates the sample support from the deferred calculation. Streaming input
continues during calculation, but those newer samples belong to a later frame.

\begin{figure}[tb]
\centering
\begin{tikzpicture}[x=0.9cm,y=0.8cm,>=Stealth,font=\small]
\draw[->] (0,0)--(12.9,0) node[right] {time};
\fill[ink!14] (0.5,0.35) rectangle (7.1,1.05);
\draw[ink] (0.5,0.35) rectangle (7.1,1.05);
\node at (3.8,0.7) {captured window: $N$ samples};
\fill[accent!18] (7.1,1.4) rectangle (11.5,2.1);
\draw[accent] (7.1,1.4) rectangle (11.5,2.1);
\node at (9.3,1.75) {$L$ batches};
\draw[dashed] (7.1,-0.15)--(7.1,2.5);
\draw[dashed] (11.5,-0.15)--(11.5,2.5);
\node[below,align=center] at (7.1,-0.15) {snapshot\\$t_r$};
\node[below,align=center] at (11.5,-0.15) {publish\\$t_r+L$};
\draw[->,ink] (7.3,0.7)--(12.2,0.7);
\node[above] at (9.6,0.7) {new input continues};
\end{tikzpicture}
\caption{Logical timeline, not to scale. The transform operates on an owned copy of the captured frame while the input ring continues to advance. A smaller batch quota delays publication without changing that frame's spectral contents.}
\label{fig:timeline}
\end{figure}


\subsection{Consequences for temporal smoothing}
The original modules use an exponential magnitude smoother. Writing its intended time
constant as $\tau$ and assuming a nominal interval $H/f_s$ gives
\begin{equation}
 y_r=\alpha y_{r-1}+(1-\alpha)|X_r|,\qquad
 \alpha=\exp[-H/(f_s\tau)].
\end{equation}
The implementation's user-facing smoothing-time parameter corresponds to
$10\tau$. If actual updates arrive every $L/f_s$, then after elapsed time $t$
the old state is attenuated approximately by
$\exp[-tH/(L\tau)]$. The effective time constant is thus
\begin{equation}
 \tau_{\mathrm{eff}}=\tau L/H.
\end{equation}
For the 4096-point example, the decay time is approximately 8.3\% shorter than
its nominal interpretation. The one-hop successor removes this completion-driven discrepancy; the Fourier module
still truncates its hop control to integer samples (Section~\ref{sec:one-hop}). The analysis also motivates timestamping
spectrogram columns by captured frames rather than assuming the configured
horizon is their spacing.

\subsection{Exact-horizon alternative}
If an application requires publication exactly every $H$ calls, it can separate
the frame clock from completion and retain a finished result until its scheduled
publication. Alternatively, for an ordered sequence of $B$ work units, assign
call $s\in\{1,\ldots,H\}$ the quota
\begin{equation}
 q_s=\floor{sB/H}-\floor{(s-1)B/H}.
 \label{eq:balanced}
\end{equation}
These quotas sum to $B$ by telescoping; each is either $\floor{B/H}$ or
$\ceil{B/H}$. For $B>0$, the last quota is positive, so the ordered work ends
on call $H$. Among schedules for $B$ unit-cost operations in $H$ calls, no
schedule can have a maximum quota smaller than $\ceil{B/H}$, by the pigeonhole
principle. Equation~\eqref{eq:balanced} attains that lower bound. This elementary
balanced schedule is not a claimed new scheduling result. The successor applies
it to a larger work sequence (Section~\ref{sec:one-hop}). Its unit-cost premise
does not make unequal operations interchangeable in execution time.

